\section{Advanced Numerical Activities}

The following activities are intended for strong master's or PhD students.
Students may use GPOPS-II, Dymos, or an equivalent direct-collocation package
where specified.

\subsection*{Problem 1: Minimum-Time Control with an Active State Constraint}

Consider
\begin{equation}
\dot{x}_1 = x_2,
\qquad
\dot{x}_2 = u,
\end{equation}
with
\begin{equation}
\mathbf{x}(0)=
\begin{bmatrix}
0\\
0
\end{bmatrix},
\qquad
\mathbf{x}(t_f)=
\begin{bmatrix}
1\\
0
\end{bmatrix},
\end{equation}
and constraints
\begin{equation}
-1 \leq u(t) \leq 1,
\qquad
|x_2(t)| \leq 0.6.
\end{equation}

Minimize
\begin{equation}
J=t_f.
\end{equation}

\begin{enumerate}
    \item Derive the optimal control structure and identify all bang and
    state-constrained arcs.
    
    \item Compute the switching times and the optimal final time analytically.
    
    \item Solve the problem using GPOPS-II or Dymos.
    
    \item Plot $x_1(t)$, $x_2(t)$, and $u(t)$, and mark every switching point.
    
    \item Repeat the numerical solution using 20, 40, and 80 mesh intervals,
    and report the error in $t_f$.
    
    \item Verify the final solution using independent high-accuracy integration.
\end{enumerate}


\subsection*{Problem 2: Plant and Actuator Co-Design of a Dynamic System}

Consider
\begin{equation}
\dot{x}_1=x_2,
\end{equation}
\begin{equation}
m\dot{x}_2=-kx_1-cx_2+u,
\end{equation}
where
\begin{equation}
m=1,
\qquad
c=0.5.
\end{equation}

The spring stiffness $k$ and actuator capacity $F_{\max}$ are plant-design
variables:
\begin{equation}
0.5 \leq k \leq 8,
\qquad
0.2 \leq F_{\max} \leq 5.
\end{equation}

The control satisfies
\begin{equation}
|u(t)| \leq F_{\max}.
\end{equation}

Use
\begin{equation}
\mathbf{x}(0)=
\begin{bmatrix}
1\\
0
\end{bmatrix},
\qquad
\mathbf{x}(t_f)=
\begin{bmatrix}
0\\
0
\end{bmatrix},
\qquad
t_f=4.
\end{equation}

Minimize
\begin{equation}
J=
\int_0^{t_f}
\left(
x_1^2+0.1x_2^2+0.01u^2
\right)\,dt
+
0.03F_{\max}^2
+
0.002k^2.
\end{equation}

\begin{enumerate}
    \item Nondimensionalize the states, control, time, and design variables.
    
    \item Formulate the simultaneous direct-collocation problem.
    
    \item Solve the problem using GPOPS-II or Dymos.
    
    \item Determine whether the actuator constraint is active.
    
    \item Repeat the solution after multiplying the $F_{\max}^2$ penalty by
    $0.1$, $1$, and $10$.
    
    \item Explain how the actuator-cost weight changes both $k^*$ and
    $F_{\max}^*$.
\end{enumerate}


\subsection*{Problem 3: Multiple Shooting versus Direct Collocation for an Unstable System}

Consider
\begin{equation}
\dot{\mathbf{x}}
=
\begin{bmatrix}
0 & 1\\
4 & 0
\end{bmatrix}
\mathbf{x}
+
\begin{bmatrix}
0\\
1
\end{bmatrix}u,
\end{equation}
with
\begin{equation}
\mathbf{x}(0)=
\begin{bmatrix}
1\\
0
\end{bmatrix},
\qquad
\mathbf{x}(3)=
\begin{bmatrix}
0\\
0
\end{bmatrix}.
\end{equation}

Minimize
\begin{equation}
J=
\int_0^3
\left(
\mathbf{x}^{T}\mathbf{x}
+
0.05u^2
\right)\,dt,
\qquad
|u|\leq 5.
\end{equation}

\begin{enumerate}
    \item Derive the exact state-transition equation over one shooting interval:
    \begin{equation}
    \mathbf{x}_{i+1}
    =
    e^{A\Delta t}\mathbf{x}_i
    +
    \int_0^{\Delta t}
    e^{A(\Delta t-\tau)}B\,u_i\,d\tau.
    \end{equation}
    
    \item Formulate a 10-segment multiple-shooting nonlinear program with
    piecewise-constant control.
    
    \item Formulate a Hermite--Simpson transcription using 30 intervals.
    
    \item Solve both formulations.
    
    \item Compare Jacobian sparsity, condition number, optimizer iterations,
    and sensitivity to initial guesses.
    
    \item Explain why direct single shooting performs poorly for this unstable
    system.
\end{enumerate}


\subsection*{Problem 4: Analytical Hermite--Simpson Jacobian}

Consider the nonlinear, parameter-dependent system
\begin{equation}
\dot{x}_1=x_2,
\end{equation}
\begin{equation}
\dot{x}_2=-p\sin x_1+u,
\end{equation}
where $p$ is a plant-design variable.

For interval $i$, define
\begin{equation}
\mathbf{f}_i
=
\mathbf{f}(\mathbf{x}_i,u_i,p),
\qquad
\mathbf{f}_{i+1}
=
\mathbf{f}(\mathbf{x}_{i+1},u_{i+1},p),
\end{equation}
and
\begin{equation}
\mathbf{x}_{i+1/2}
=
\frac{\mathbf{x}_i+\mathbf{x}_{i+1}}{2}
+
\frac{h_i}{8}
\left(
\mathbf{f}_i-\mathbf{f}_{i+1}
\right).
\end{equation}

The Hermite--Simpson defect is
\begin{equation}
\mathbf{d}_i
=
\mathbf{x}_{i+1}
-
\mathbf{x}_i
-
\frac{h_i}{6}
\left(
\mathbf{f}_i
+
4\mathbf{f}_{i+1/2}
+
\mathbf{f}_{i+1}
\right).
\end{equation}

\begin{enumerate}
    \item Derive
    \begin{equation}
    \frac{\partial\mathbf{d}_i}{\partial\mathbf{x}_i},
    \qquad
    \frac{\partial\mathbf{d}_i}{\partial\mathbf{x}_{i+1}},
    \qquad
    \frac{\partial\mathbf{d}_i}{\partial u_i},
    \qquad
    \frac{\partial\mathbf{d}_i}{\partial u_{i+1}},
    \qquad
    \frac{\partial\mathbf{d}_i}{\partial p}.
    \end{equation}
    
    \item Assemble the Jacobian sparsity pattern for $N$ intervals.
    
    \item Determine the exact number of nonzero entries.
    
    \item Explain why the column associated with $p$ is globally dense while
    the state and control blocks remain locally sparse.
    
    \item Verify the analytical Jacobian using complex-step differentiation.
\end{enumerate}


\subsection*{Problem 5: Derivative Verification for a Shooting Formulation}

Consider
\begin{equation}
\dot{x}=-x^3+u,
\qquad
x(0)=1,
\qquad
0\leq t\leq 2.
\end{equation}

Parameterize $u(t)$ using $N=20$ piecewise-constant values:
\begin{equation}
\mathbf{U}
=
\begin{bmatrix}
u_0 & u_1 & \cdots & u_{19}
\end{bmatrix}^{T}.
\end{equation}

Minimize
\begin{equation}
J(\mathbf{U})
=
x(2)^2
+
0.05
\int_0^2u(t)^2\,dt.
\end{equation}

\begin{enumerate}
    \item Derive the forward sensitivity equations
    \begin{equation}
    \frac{d}{dt}
    \left(
    \frac{\partial x}{\partial u_j}
    \right)
    =
    -3x^2
    \frac{\partial x}{\partial u_j}
    +
    \frac{\partial u}{\partial u_j}.
    \end{equation}
    
    \item Derive the gradient $\nabla_{\mathbf{U}}J$.
    
    \item Compute the gradient using:
    \begin{enumerate}
        \item forward finite differences,
        \item central finite differences,
        \item complex-step differentiation,
        \item sensitivity equations or automatic differentiation.
    \end{enumerate}
    
    \item Perform a directional-derivative test using a random normalized vector
    $\mathbf{p}$.
    
    \item Plot derivative error for step sizes from $10^{-2}$ to $10^{-14}$.
    
    \item Explain the truncation-error and cancellation-error regions in the
    resulting plot.
\end{enumerate}


\subsection*{Problem 6: Hypersensitive Optimal-Control Problem}

Consider
\begin{equation}
\dot{x}=-x^3+u,
\end{equation}
with
\begin{equation}
x(0)=1,
\qquad
x(T)=1.5,
\qquad
T=50.
\end{equation}

Minimize
\begin{equation}
J=
\frac{1}{2}
\int_0^T
\left(
x^2+u^2
\right)\,dt.
\end{equation}

\begin{enumerate}
    \item Solve the problem using direct single shooting.
    
    \item Solve it using multiple shooting with 5, 10, and 20 segments.
    
    \item Solve it using GPOPS-II or Dymos with an adaptive collocation mesh.
    
    \item Plot the state and control near:
    \begin{enumerate}
        \item the initial boundary layer,
        \item the long interior arc,
        \item the terminal boundary layer.
    \end{enumerate}
    
    \item Compare the methods in terms of:
    \begin{enumerate}
        \item convergence,
        \item sensitivity to the initial guess,
        \item Jacobian conditioning,
        \item mesh size,
        \item computation time.
    \end{enumerate}
    
    \item Independently integrate the collocation solution and report
    \begin{equation}
    e_x=
    \max_t
    \left|
    x_{\mathrm{collocation}}(t)
    -
    x_{\mathrm{verification}}(t)
    \right|.
    \end{equation}
    
    \item Explain why this problem is difficult for single shooting even though
    it contains only one state.
\end{enumerate}
